The goal of this section is to prove Theorem~\ref{thm:cleaning-intro}.

\cleaning*

Towards the proof of Theorem~\ref{thm:cleaning-intro}, we first prove Lemma~\ref{lem:cleaning-wrap-up}. 

\begin{lemma}[Identifying to a \hardpattern{t}]\label{lem:cleaning-wrap-up}
Let $\mathcal{D}$ be a class of graphs that is closed under transitive equivalence and identifying vertices. Let $D \in \mathcal{D}$ and let $t$ be a positive integer. 
Then there exists $t'$ that depends only on $t$,
such that if $D$ contains a $t'$-\toughpair\ then there exists $D' \in \mathcal{D}$ that is a \hardpattern{t}.
\end{lemma}

From Lemma~\ref{lem:cleaning-wrap-up}, the proof of Theorem~\ref{thm:cleaning-intro} follows easily. We give this proof before proving Lemma~\ref{lem:cleaning-wrap-up}.

\begin{proof}[Proof of Theorem~\ref{thm:cleaning-intro}]

We first prove the forward direction.
From Lemma~\ref{lem:cleaning-wrap-up}, for every positive integer $t$, $\mathcal{D}$ contains a digraph from some \hardpattern{t}, that is, $\mathcal{D}$ contains a digraph from some $\mathcal{C}_i$. Since $\kappa$ is finite, 
there exists $i \in [\kappa]$ such that infinitely many 
digraphs of $\mathcal{C}_i$ belong to $\mathcal{D}$. Further, since $\mathcal{D}$ is closed under vertex identifications and from any member of $\mathcal{C}_i$ every smaller member can be obtained by vertex identifications, we conclude that $\mathcal{C}_i \subseteq \mathcal{D}$.

For the backward direction observe from the definitions of $\{ \mathcal{C}_1, \ldots ,\mathcal{C}_{\kappa}\}$, that each $t$-hard digraph of each $\mathcal{C}_i$ contains in fact a $t$-\toughpair.
\end{proof}

The remainder of this section is dedicated to the proof of Lemma~\ref{lem:cleaning-wrap-up} which is divided into five separate steps presented in Sections~\ref{sec:firstramsey}, \ref{sec:simplifyingbiclique}, \ref{sec:semicleaning}, \ref{sec:semicleantoclean} and \ref{sec:cleaningbiclique}, respectively. Below we give the main results of each of these sections and show how they together imply Lemma~\ref{lem:cleaning-wrap-up}. Before stating the results, we give some important definitions that are used throughout the section.

\paragraph*{Definitions.}

 For any digraph $D$, we denote by $D^{\star}$ the {\em transitive closure} of $D$, that is $D^{\star}$ is obtained from $D$ by repeatedly adding edges $(u,v)$ whenever $(u,v)$ is not already an edge but there is a $(u,v)$-path in $D$.

For a digraph $D$ and two ordered sets $A =(a_1,\ldots, a_t),B=(b_1, \ldots,b_t) \subseteq V(D)$ such that $|A|=|B|$, we say that $D$ has an {\em $(A,B)$-matching} if for each $i \in [t]$, $(a_i,b_i) \in E(D)$. We say that $D$ has an {\em $(A,B)$-induced-matching} if $D$ has an $(A,B)$-matching, for each $i \neq j$, $(a_i,b_j) \not \in E(D)$, $(b_j,a_i) \not \in E(D)$, $(b_i,a_i) \not \in E(D)$ and $A,B$  are independent sets in $D$.

We say that $D$ has an {\em $(A,B)$-biclique} if for each $i,j \in [t]$, $(a_i,b_j) \in E(D)$ and $(b_j,a_i) \not \in E(D)$. We say that $D$ has an {\em $(A,B)$-induced-biclique}
if $D$ has an $(A,B)$-biclique and, $A, B$ are independent sets in $D$. We say that $D$ has a {\em $t$-biclique} if there exists an $(A,B)$-biclique in $D$ for some $A,B \subseteq V(D)$ and $|A|=|B|=t$. Similarly, we say that $D$ has a {\em $t$-induced-biclique} if there exists an $(A,B)$-induced-biclique in $D$ for some $A,B \subseteq V(D)$ and $|A|=|B|=t$. In all these cases, we call the set of edges $\{(a_i,b_i): i \in [t]\}$ as the {\em matching edges} of $(A,B)$.

For positive integers $i,j$, let $\mathcal{R}(i,j)$ denote the minimum number of vertices such that any $\mathcal{R}(i,j)$-vertex complete graph whose edges are colored with $j$ colors, contains a monochromatic clique of size $i$. Let $\mathcal{R}'(i)$ denote the minimum integer such that any bipartite graph that has a matching of size $\mathcal{R}'(i)$, either has an induced matching of size $t$ or a $t$-induced-biclique.
From Ramsey Theorem, such numbers always exists and they depend only on $i,j$, or $i$, respectively.
Recall the definitions of weakly independence, strongly independence and $t$-\toughpair\ from Section~\ref{sec:prelims}.

\begin{definition}[Ordered $t$-\toughpair]
Given a 
digraph $D$, $E_1, E_2 \subseteq E(D^{\star})$, we say that $(E_1,E_2)$ is an ordered $t$-\toughpair\  in $D$ if 
\begin{enumerate}
\item $|E_1| = |E_2|=t$,
\item all edges in $E_i$ are pairwise weakly independent in $D$, for each $i \in \{1,2\}$, and
\item for each $e_1 \in E_1$ and $e_2 \in E_2$, $(e_1,e_2)$ are strongly independent in $D$, and
\end{enumerate}
there exists an ordering of the sets $\head(E_1) = (w_1, \ldots, w_t)$, $\tail(E_1)=(x_1, \ldots, x_t)$, $\head(E_2)=(y_1, \ldots, y_t)$ and $\tail(E_2)=(z_1, \ldots, z_t)$ such that
\begin{enumerate}
\item for any $1\leq i < j \leq t$, there is no $(w_j,x_i)$-path, no $(w_j,w_i)$-path and no $(x_j,x_i)$-path in $D$,
\item for any $1\leq i < j \leq t$, there is no $(y_j,z_i)$-path, no $(y_j,y_i)$-path and no $(z_j,z_i)$-path in $D$.
\end{enumerate}
If $(E_1,E_2)$ is an ordered $t$-\toughpair\ then we treat $E_1$ and $E_2$ as ordered sets such that their head sets and tail sets satisfies the above properties. We say ordered \toughpair\ to mean an ordered $t$-\toughpair\ for some $t$.
\end{definition}

Observe that, unlike the $t$-\toughpair, if $(E_1,E_2)$ is an ordered $t$-\toughpair\ in $D$, then the edges of $E_1$ and $E_2$ may not be minimal in $D$.

\paragraph*{Simplifying the $t$-\toughpair.}
In Section~\ref{sec:firstramsey}, using Ramsey arguments we show that it $D$ contains a $t$-\toughpair\ for a large enough $t$, then it contains one of the three structures described in Lemma~\ref{lem:pre-cleaning}.

\begin{restatable}[Hard sub-structures]{lemma}{firstramsey}\label{lem:pre-cleaning}
Let $D$ be a reachability-minimal digraph and $t$ be a positive integer. 
If $D$ contains an $\mathcal{R}(2t,9)$-\toughpair\
then one of the following holds.
\begin{enumerate}
\item $D$ contains an ordered $t$-\toughpair\, or
\item $D$ contains a $t$-biclique, or
\item there exist ordered sets $A,B \subseteq V(D)$ such that $|A|=|B|=t$, $D$ contains an $(A,B)$-matching and $D^{\star}$ contains an $(A,B)$-induced-biclique. 
\end{enumerate}
\end{restatable}

Note that if $D$ is reachability-minimal and contains a $t$-biclique then it contains a $t$-induced-biclique. 
From Lemma~\ref{lem:pre-cleaning} one concludes that, in order to prove Lemma~\ref{lem:cleaning-wrap-up}, it is enough to identify one of the three structures defined in Lemma~\ref{lem:pre-cleaning} to some \hardpattern{t}. The first outcome of Lemma~\ref{lem:pre-cleaning} is handled in Section~\ref{sec:cleaningordered}, the second outcome is handled in Section~\ref{sec:cleaningbiclique} and the third outcome is handled in Section~\ref{sec:simplifyingbiclique}.

\paragraph*{Cleaning ordered $t$-\toughpair.} In Section~\ref{sec:cleaningordered} we show that if the outcome of Lemma~\ref{lem:pre-cleaning} is an ordered $t$-\toughpair, then one can identify $D$ to either a \hardpattern{t'} or $D^{\star}$ contains a $t'$-induced-biclique whose edges are minimal in $D$. 
This is achieved by two rounds of cleaning: in the first round the digraph is cleaned to an intermediate structure, called a {\em semi-cleaned ordered $t$-\toughpair} (defined below), and in the second round this semi-cleaned ordered $t$-\toughpair\ is either identified to obtain a \hardpattern{t'}, or it can be shown that it contains a $t'$-induced-biclique whose edges are minimal.
This is formalized in Lemmas~\ref{lem:semi-cleaning} and~\ref{lem:semicleantoclean} which are proved in Section~\ref{sec:semicleaning} and~\ref{sec:semicleaning}, respectively.

\begin{definition}[Semi-cleaned ordered $t$-\toughpair]
We say that a digraph $D$ is a semi-cleaned ordered $t$-\toughpair, if it contains an ordered $t$-\toughpair\ $(A,B)$ such that the vertex set of $D$ contains at most two vertices, called a source $\sss$ and a sink $\ttt$, outside of the set $V(A \cup B)$, and $N^-_D(\sss), N^+_D(\ttt) =\emptyset$. (Note that from this definition $(\ttt,\sss) \not \in E(D)$).
\end{definition}

\begin{restatable}[Semi-cleaning the ordered \toughpair]{lemma}{semicleaning}\label{lem:semi-cleaning}
Let $D$ be a
digraph such that $D$ contains an ordered $t^2$-\toughpair, then one can obtain $\widehat{D}$ from $D$ by identification, such that $\widehat{D}$ is a semi-cleaned $t$-ordered \toughpair, for some function $g$ that depends only on $t$.
\end{restatable}

For any positive integer $t$, let $\htt=2\mathcal{R}(2\mathcal{R}(2\mathcal{R}(4t+2,4),4),5)$.

\begin{restatable}[Cleaning the semi-cleaned ordered \toughpair]{lemma}{semicleantoclean}
\label{lem:semicleantoclean}
If $D$ is a semi-cleaned ordered $\htt$-\toughpair\ 
then either, 
\begin{itemize}
\item $D^{\star}$ contains a $t$-induced-biclique whose edges are minimal in $D$, or
\item $D$ can be identified to a digraph $\hat{D}$ such that $\hat{D}$ is transitively equivalent to a \hardmatching{t}.
\end{itemize}
\end{restatable}

For any positive integer $t$, let $\cor{t}= (\htt)^2$.
Corollary~\ref{cor:cleaning-ordered} follows from Lemmas~\ref{lem:semi-cleaning} and~\ref{lem:semicleantoclean}.

\begin{corollary}\label{cor:cleaning-ordered}
If a digraph $D$ contains an ordered $\cor{t}$-\toughpair,
then either $D$ can be identified to a \hardmatching{t}, or $D^{\star}$ contains a $t$-induced-biclique whose edges are minimal in $D$.
\end{corollary}

If the outcome of Corollary~\ref{cor:cleaning-ordered} a \hardmatching{t}, then we are done. Otherwise, we need to clean the $t$-induced-biclique whose edges are minimal,
which is what is done next.
\paragraph*{Cleaning a minimal induced-biclique.} In Section~\ref{sec:cleaningbiclique} we show that if $D$ contains a $9t$-induced-biclique whose edges are minimal, then $D$ can be identified to a \hardbiclique{t}. In this case, Lemma~\ref{lem:cleaning-wrap-up} is proved.

\begin{restatable}[Cleaning minimal biclique]{lemma}{cleaningbiclique}\label{lem:cleaningbiclique}
For a positive integer $t$ and a digraph $D$, if $D^{\star}$ contains a $9t$-induced-biclique whose edges are minimal in $D$, then $D$ can be identified to digraph that is transitively equivalent to a \hardbiclique{t}.
\end{restatable}

\paragraph*{Simplifying the third outcome of Lemma~\ref{lem:pre-cleaning}.} In Section~\ref{sec:simplifyingbiclique}, we show that if Lemma~\ref{lem:pre-cleaning} outputs its third outcome, then one can contract some edges of the input digraph (and hence obtain a digraph in the same pattern class as the input, since the pattern class is closed under identification and contraction is a type of identification) such that the resulting digraph either contains an ordered $t'$-\toughpair\ (which we know how to deal using Corollary~\ref{cor:cleaning-ordered}), or it contains a $t'$-biclique whose edges are minimal (which we also know to deal with because of Lemma~\ref{lem:cleaningbiclique}).

For positive integers $t,p$, define $f(t,p) = 2^{\mathcal{R}'(2t)} \cdot f(t,p-1) + \mathcal{R}'(2t)$ when $p >1$ and $f(t,1)=t$.

\begin{restatable}[Simplifying a biclique]{lemma}{digging}\label{lem:digging}
Let $D$ be a directed graph and $t$ be a positive integer.  
Suppose there exists ordered sets
$A,B \subseteq V(D)$ such that $|A|=|B|\geq f(t,2t+2)$, 
$D^{\star}$ contains an $(A,B)$-induced-biclique and the matching edges of $(A,B)$ are minimal in $D$,.
Then, one can obtain $\widehat{D}$ from $D$ by contraction, such that 
either
\begin{enumerate}
\item $\widehat{D}$ contains an ordered $t$-\toughpair, or
\item $\widehat{D}^*$ contains a $t$-induced-biclique whose edges are minimal in $\widehat{D}$.
\end{enumerate}
\end{restatable}

\begin{proof}[Proof of Lemma~\ref{lem:cleaning-wrap-up}]
We now give a proof of Lemma~\ref{lem:cleaning-wrap-up} using the lemmas stated earlier. The proofs of the those lemmas appear in their respective sections.

Set $t_1 = \max\{\cor{9t},9t\}$, $t_2 =\max \{ f(t_1, 2t_1+2), \cor{9t_1}, 9t_1 \}$ where $f$ is defined as in Lemma~\ref{lem:digging} and $\corr$ is defined in Corollary~\ref{cor:cleaning-ordered}. Then, define $t' =  \mathcal{R}(2t_2,9)$. We will now show that if $D$ contains a $t'$-\toughpair, then there exists $D' \in \mathcal{D}$ that is a \hardpattern{t}.

Given $D$, let $(X,Y)$ be a $t'$-\toughpair\ in $D$. Let $D_1$ be a spanning subgraph of $D$ such that $D_1$ is reachability-minimal, $D_1$ is transitively equivalent to $D$ and $(X,Y)$ is a $t'$-\toughpair\ in $D_1$. Observe that such a graph $D_1$ exists and can be obtained by starting with the edge set $E_0=X \cup Y$ and adding an edge of $D$ to $E_0$ as long as $D[E_0]$ remains reachability-minimal. 

Using Lemma~\ref{lem:pre-cleaning} on $D_1$, we conclude that either $D_1$ contains an ordered $t_2$-\toughpair, or a $t_2$-biclique, or there exists $A,B \subseteq V(D_1)$ such that $|A|=|B|=t_2$ and $D_1$ contains an $(A,B)$-induced-matching and $D_1^{\star}$ contains an $(A,B)$-induced-biclique. 

In the first case, by applying Corollary~\ref{cor:cleaning-ordered} on $D^{\star}_1$, we either get $D_2 \in \mathcal{D}$ such that $D_2$ is some \hardmatching{t}, 
or conclude that $D^{\star}_1$ contains a $9t$-induced-biclique whose edges are minimal in $D_1$. Then applying Lemma~\ref{lem:cleaningbiclique} to $D_1$ gives some \hardbiclique{t} in $\mathcal{D}$.
In the second case, again by applying Lemma~\ref{lem:cleaningbiclique} to $D_1$ gives some \hardbiclique{t} in~$\mathcal{D}$.

In the third case, by applying Lemma~\ref{lem:digging} we conclude that there exists $D_2 \in \mathcal{D}$ such that either $D_2$ contains an ordered $t_1$-\toughpair, in which case a further application of Corollary~\ref{cor:cleaning-ordered} on $D^{\star}_2$ either yields a $D_3 \in \mathcal{D}$ which is a \hardmatching{t}, or we conclude that $D^{\star}_2$ contains a $9t$-induced-biclique whose edges are minimal in $D_2$. In the remaining cases, applying Lemma~\ref{lem:cleaningbiclique} yields a $D_4 \in \mathcal{D}$ which is some \hardbiclique{t}.
\end{proof}

\paragraph*{Basic terminology for the remaining section.}Throughout the remaining section, we use the following basic notation. For two integers $i,j$, $[i]$ denotes the set $\{1, \ldots, i\}$ and $[i,j]$ denotes the set $\{i,i+1, \ldots,j\}$.  For any (di)graph $D$ and sets $S,T \subseteq V(D)$, an $(S,T)$-path in $D$ is a path in $D$ from some vertex of $S$ to some vertex of $T$. If $S$ or $T$ is singleton, say $S=\{v\}$, then we use the notation $(v,T)$-path. Let $e \in E(D)$, let $D/e$ represents the (di)graph obtained after contracting $e$, that is, by deleting the endpoints of $e$ and adding a new vertex $x_e$ such that the set of in-neighbours (resp.~out-neighbours) of $x_e$ is the union of the set of in-neighbours (resp.~out-neighbours) of the end points of $e$ in $D$. For $E' \subseteq E(D)$, $V(E') \subseteq V(G)$ is the set of vertices that are endpoints of some edge in $E'$. If $D$ is a digraph, the for any $v \in V(D)$, $N^-_D(v)$ denotes the set of in-neighbours if $v$ in $D$, $N^+_D(v)$ denotes the set of out-neighbours if $v$ in $D$ and $N_D(v)$ denotes the set of in-neighbours and out-neighbours, called the neighbours, of $v$ in $D$.
We say that a digraph $D$ is connected if its underlying undirected graph is connected. 